\chapter{FFI Foundations and Semantic Contracts}
\label{ch:ffi-foundations}

A native function can receive the right bits and still be called incorrectly.
Its pointer may refer to an object the collector can move, its length may count characters instead of bytes, or its callback may reenter a runtime that the caller assumed was suspended.
Machine-level agreement is necessary, but the whole boundary must preserve the meanings of values, effects, and lifetimes.
A \term{foreign function interface}, or FFI, is the machinery and contract that lets one execution environment invoke operations defined by another.

The running example exposes checked integer addition, owned byte buffers, and synchronous callbacks through a C interface.
Lua and CPython wrappers share that interface while retaining their own language-level error and ownership rules.
Chapter~\ref{ch:ffi-case-study} develops the implementation; the intervening chapters derive its contracts and examine broader extension models.

\section{Separate three kinds of agreement}
\label{sec:ch69-layers}
Suppose a function adds two signed 32-bit integers and writes its result through an output pointer.
Its source declaration determines parameter types, but does not by itself say whether overflow wraps, traps, or produces a status code.
Our interface returns \code{BB\_RANGE} on overflow and leaves the output untouched.
The caller must also supply valid writable storage for that output during the call.

The ABI decides where the two integers and pointer travel and where the status returns.
The adapter decides how a Python or Lua integer is converted and how a failure becomes a language error.
The semantic contract decides which inputs and effects are permitted.
Figure~\ref{fig:ch69-layers} makes these layers explicit.
Changing one layer without updating the others can produce a binary that links correctly and fails only on unusual inputs.

\begin{figure}[htbp]
\centering
\begin{tikzpicture}[node distance=11mm]
\node[basic block] (language) {language values and error rules};
\node[basic block,below=of language] (contract) {conversion, ownership, effects};
\node[basic block,below=of contract] (abi) {ABI locations and native call};
\draw[cfg edge] (language)--(contract); \draw[cfg edge] (contract)--(abi);
\end{tikzpicture}
\caption{The adapter preserves language meaning while the calling convention transports values.}
\label{fig:ch69-layers}
\end{figure}

Table~\ref{tab:ch69-agreement} gives independently failing agreements.
\term{Interoperability} therefore describes a stated combination of languages, runtimes, target ABI, data conventions, and supported operations.
It is not established merely by exporting a C symbol.

\begin{table}[htbp]
\caption{Successful symbol lookup establishes none of these additional contracts.}
\label{tab:ch69-agreement}
\centering
\begin{tabular}{ll}
\toprule
Agreement & Example mismatch \\
\midrule
Value meaning & signed integer interpreted as unsigned \\
Representation & UTF-8 byte count treated as character count \\
Calling convention & floating argument placed in an integer register \\
Lifetime & retained pointer outlives its borrowed buffer \\
Control & exception crosses an unsupported foreign frame \\
Concurrency & callback runs on an unattached runtime thread \\
\bottomrule
\end{tabular}
\end{table}

\section{Make conversions partial and explicit}
\label{sec:ch69-conversion}
A conversion from arbitrary-precision integers to \code{int32\_t} is a partial function unless the interface deliberately specifies truncation.
The accepted range is $-2^{31}$ through $2^{31}-1$.
Our Python wrapper additionally accepts exact integers only, rejecting Boolean values and integer subclasses; the Lua wrapper requires the integer subtype and rejects \code{1.0} even though it denotes an integral numeric value.
These are explicit binding policies, not universal recommendations for every FFI.

For $a=19$ and $b=23$, conversion succeeds and checked addition produces $42$.
For $a=2^{31}-1$ and $b=1$, both inputs fit, but the result does not.
The core computes in a wider integer type, detects the range failure, and does not overwrite the caller's output sentinel.
For a Python input $2^{80}$, conversion fails before native execution.
Table~\ref{tab:ch69-add} separates those two failure stages.

\begin{table}[htbp]
\caption{Input conversion and native arithmetic enforce distinct range obligations.}
\label{tab:ch69-add}
\centering
\begin{tabular}{lll}
\toprule
Inputs & Stage reached & Outcome \\
\midrule
$19,23$ & native addition & result $42$ \\
$2^{31}-1,1$ & native addition & range status; output unchanged \\
$2^{80},0$ & adapter conversion & rejected before native call \\
Boolean and integer & adapter type check & rejected by this binding \\
\bottomrule
\end{tabular}
\end{table}

An adapter specification describes argument direction, nullability, extent units, ownership transfer, and failure behavior in addition to scalar types.
A \code{const} pointer says this particular C access path does not write through it; another alias may still mutate the storage.
A pointer to a constant object is also different from a pointer whose target remains alive indefinitely.
Algorithm~\ref{alg:ch69-admit} checks a declaration's completeness before generating a wrapper.

\begin{algorithm}[htbp]
\caption{Admit a foreign declaration with a complete supported boundary contract.}
\label{alg:ch69-admit}
\begin{pseudocode}
\Require declaration, target ABI, supported conversion and lifetime policies
\Ensure wrapper plan or explicit unsupported-contract diagnostic
\State validate parameter/result types and calling convention
\ForAll{pointer-like parameters}
  \State require extent units, null policy, direction, and lifetime policy
\EndFor
\State identify callbacks, reentrancy, thread restrictions, and retained references
\State identify the error channel and output validity on failure
\State reject requirements without an implemented conversion or cleanup path
\State return the admitted plan and its assumptions
\end{pseudocode}
\end{algorithm}

For $A$ parameters and $K$ explicit contract annotations, structural admission costs $O(A+K)$ plus target-specific type processing.
It does not prove that the library obeys its declaration.
That trust boundary requires compatible headers, implementation review, and tests; guessing a signature from a symbol name cannot establish it.

\section{Describe foreign effects to the optimizer}
\label{sec:ch69-effects}
The compiler must know more than the return type to optimize around a foreign call.
A call can modify reachable memory, retain a pointer, invoke a callback, block, or fail to return.
An imported effect summary must describe all allowed behavior, including behavior reached indirectly through callbacks.

Consider loading a global $g$, calling a library visitor, and loading $g$ again.
The visitor's own code may contain no store to $g$, yet its callback may change $g$.
Replacing the second load with the first requires an effect guarantee covering the callback path.
Likewise, a no-capture property must specify whether the library may retain the pointer after a callback or exceptional return.
Algorithm~\ref{alg:ch69-summary} gives conservative treatment of such calls.

\begin{algorithm}[htbp]
\caption{Apply an imported effect contract conservatively.}
\label{alg:ch69-summary}
\begin{pseudocode}
\Require call, current analysis facts, trusted or conservatively inferred summary
\Ensure retained facts remain valid after every allowed behavior
\State expand the summary to include permitted callback effects
\State invalidate memory facts overlapping possible writes
\State mark retained references as escaped under the ownership policy
\State preserve required exceptional and nonreturning control paths
\State keep only facts justified by the complete summary
\end{pseudocode}
\end{algorithm}

With $F$ tracked facts and an overlap query costing $Q$, a direct invalidation scan costs $O(FQ)$; indexed summaries can avoid scanning unrelated regions.
The annotation itself is not a proof.
An incorrect read-only declaration can miscompile a program even when its register transfers are perfect.

\section{Specify the temporal boundary}
\label{sec:ch69-time}
A borrowed byte slice can be valid for one synchronous call without being valid for a registration that returns immediately and invokes a callback later.
The contract must name the event ending the borrow.
For asynchronous work, that event may be completion acknowledgement rather than cancellation request.
Figure~\ref{fig:ch69-borrow} places a synchronous borrow between acquisition and release.

\begin{figure}[htbp]
\centering
\begin{tikzpicture}[node distance=11mm]
\node[basic block] (acquire) {acquire stable readable slice};
\node[basic block,below=of acquire] (call) {native call and allowed callbacks};
\node[basic block,below=of call] (release) {native use ends; release slice};
\draw[cfg edge] (acquire)--(call); \draw[cfg edge] (call)--(release);
\end{tikzpicture}
\caption{The borrow includes all allowed activity during the call, including reentrancy.}
\label{fig:ch69-borrow}
\end{figure}

The implementation avoids a retained borrow by copying byte inputs into a native allocation.
That costs linear time and storage in the byte count but makes later ownership explicit.
Zero-copy alternatives in Chapter~\ref{ch:ffi-marshaling} require stronger lifetime and mutation guarantees.

\section{Exercises}
\label{sec:ch69-exercises}
\begin{enumerate}
\item\label{ex:ch69-range} Work the three numeric rows in Table~\ref{tab:ch69-add} with output initially $99$.
Identify which failures enter native code and the output after each call.
\item\label{ex:ch69-extent} Use ASCII A followed by U+00E9 to show why a UTF-8 pointer's length must specify its unit.
\item\label{ex:ch69-effect} A visitor's callback changes global $g$ from $4$ to $9$.
Show why reusing a pre-call load after the visitor is unsound.
\item\label{ex:ch69-retain} A library retains a pointer after returning.
List the lifetime and completion facts needed to accept a borrowed-buffer implementation.
\end{enumerate}

\section{Further reading}
\label{sec:ch69-reading}
C11 specifies scalar, pointer, and object-lifetime rules at the native side~\cite{wg14-2011-c11}.
The Lua and Python manuals document concrete runtime conversion and error mechanisms~\cite{ierusalimschy-2026-lua,python-2026-capi}.
The libffi manual explains how a known signature becomes an executable call while leaving ownership and higher-level semantics to its caller~\cite{libffi-2026-manual}.
